\chapter{Durchführung} 
\label{ch:Durchfuehrung}

\section{Messverfahren}
\label{sec:Messverfahren}
Die experimentelle Durchführung gliedert sich in mehrere Teilmessungen zur Untersuchung der verschiedenen Strahlungsarten und ihrer Wechselwirkungen mit Materie. Zu Beginn jeder Messreihe wird der Nulleffekt über einen definierten Zeitraum ohne radioaktive Quelle bestimmt, um diesen von den späteren Messwerten zu subtrahieren. Für die Untersuchung von $\beta$-Strahlung wird ein Strontium-Präparat verwendet, wobei sukzessive Aluminium-Absorberplatten bekannter Dicke zwischen Quelle und Zählrohr platziert werden, bis die Zählrate auf den Wert des Nulleffekts abgefallen ist. Analog dazu erfolgt die Untersuchung von $\gamma$-Strahlung unter Verwendung eines Kobalt-Präparats, bei dem Bleiplatten als Absorbermaterial schichtweise eingebracht werden, um das Lambert-Beersche Schwächungsverhalten zu verifizieren. Zur Bestimmung der Aktivität des $\gamma$-Strahlers wird die Zählrate bei variierenden Abständen zwischen Detektor und Quelle aufgenommen. Für die $\alpha$-Strahlung wird ein Amerizium-Präparat in einem evakuierbaren Glaszylinder genutzt. Durch kontrollierte Variation des Luftdrucks im Zylinder lässt sich die effektive Reichweite der $\alpha$-Teilchen präzise bestimmen, da sich hierbei die Dichte des Gasvolumens zwischen Quelle und Zählrohr stetig verändert.

\iimg[0.61]{plt1}{plt2}{(a) Zeigt die Reichweite von $\beta$-Strahlung in Aluminum und von $\alpha$-Strahlung in Luft. (b) zeigt den auf die Masse normierten Wirkungsquerschnitt von Photonen in Blei und Aluminum in der Abhänigkeit der Photonenenergie. \cite{skriptPAP22}}

\section{Versuchsaufbau}
\label{sec:Versucshaufbau}
Der Versuchsaufbau besteht im Wesentlichen aus einer Strahlungsquelle, einem variablen Absorbersystem sowie einem Detektor zur Registrierung der ionisierenden Strahlung. Als Detektor dient ein Geiger-Müller-Zählrohr, welches über ein Zählgerät zur Registrierung und Summierung der elektrischen Impulssignale angeschlossen ist. Für die Messung der $\alpha$-Strahlung ist das Geiger-Müller-Zählrohr gemeinsam mit dem Amerizium-Präparat in einem evakuierbaren Glaszylinder untergebracht, welcher mit einer Druckregulierung und einer Manometeranzeige versehen ist. Die Präparate für $\beta$- und $\gamma$-Messungen werden in definierten Halterungen positioniert, die eine exakte Ausrichtung und reproduzierbare Abstände zum Zählrohrfenster gewährleisten. Zwischen Quelle und Detektor können passgenaue Absorberplatten aus Aluminium beziehungsweise Blei eingepasst werden.

\subsubsection*{Versuchsaufbaus}
\label{sec:Skizze}

\img{VAufbau}{Fotographie des Versuchaufbaues. \cite{skriptPAP22}}